\begin{definition}[Hypotheses and the composite term along $\Delta$]
\label{def:c10-hyp}
The hypotheses along $\Delta$ are:
\begin{itemize}
\item (i) $\hbar F_1^{(0)}$ solves the arity-zero part of $s(\hbar F_1)=\hbar\,\mathrm{Ob}$ and contains the finite counterterm of Theorem \ref{thm:c4-finite}, and $\hbar F_1^{(1)}(\Delta)$ solves \eqref{eq:c9-arity1}, whose right side is the sum of the coefficients of $\varepsilon^0$ of the defect cycles along $\Delta$ (Proposition \ref{prop:c9-structure} (ii)), equal to $4$ times the coefficient of $\log\varepsilon$ in $I^{CT}[\varepsilon]$ plus the composite part (Proposition \ref{prop:c9-structure} (iii)); so $\hbar F_1^{(1)}(\Delta)$ minus $4$ units times the field-counting functional has the composite part as its image under $Q+\{I,-\}$. Its quadratic part is \eqref{eq:c9-F1Delta}: the translation-invariant solution of Theorem \ref{thm:c9-quadratic} (iii), whose part beyond $4$ units times the field-counting functional is $\big[-\tfrac23\int(A,A^*)+\tfrac13\int(\lambda^*,\lambda)+\tfrac43\int(c^*,c)\big]$ units, plus the $Q$-closed directions: the three sector parts of $S^{(1)}(\Delta)$, with the coefficients $\theta$, $\theta_\lambda$, $\theta_V$ units, $\theta$ the coefficient of \eqref{eq:c9-Fperp}, and $Q\int(\iota_xA|_V,c^*)$, $Q\int(\rho(x)\lambda^*,\lambda^*)$. The arity-two equations along $(\Delta,P_b)$ and $(\Delta,\xi)$ fix $\theta=\theta_\lambda=\theta_V=-\tfrac13$ (Theorem \ref{thm:c9-theta2}); the statements below are written for general values, which is how the free directions enter.
\item (ii) $O$ is homogeneous, $(Q+\{I,-\})O=0$, and $L_1(O)^{(0)}=0$ for a finite part $\psi$ of the operators that is linear in $\omega$, translation invariant and preserves the Weyl weight; such a $\psi$ removes the outputs of Remark \ref{rem:c5-cocycle}.
\item (iii) The local representative $L(O)^{(1)}(\Delta)$ exists (Definition \ref{def:c8-rep}); given (i), this holds for every local functional $O$ (Theorem \ref{thm:c11-classes} (v)).
\end{itemize}
The composite graphs of $O$ are the graphs of the linearised flow \eqref{eq:c9-D1a} in which the composite kernel $2(P\star v.P\star\otimes1)K_0$, a line from $c$ to $A|_V$ (Proposition \ref{prop:c9-composite} (ii)), closes two legs of the tree-level flow of $O$: either the cycle passes through two letters $X_a$, $X_b$ of $O$, the composite cycles through the pair, or it closes on the tree attached to one letter, which gives the renormalised composite two-point functional of Remark \ref{rem:c9-Pi} on a leg of $O$. Their singular parts are the arity-one counterterm of $O$ along $\Delta$ and the quadratic part of $I^{CT}_{\mathfrak{psl}(4|4)}[\varepsilon]^{(1)}(\Delta)$ (Proposition \ref{prop:c9-structure} (i)). The composite term of $O$ is the local functional whose tree-level flow at scale $T$ is
\[
\lim_{\varepsilon\to0}Q\big(\text{composite graphs of }O\text{ minus their counterterm trees}\big)[T]-\big(\text{tree insertion of the composite part into }O[T]\big),
\]
and the composite term of the pair $(X_a,X_b)$ is the image under $Q$ of the renormalised composite cycles through the pair.
\end{definition}
